\documentclass[11pt]{article}
\usepackage[margin=1in]{geometry}
\usepackage{amsmath,amssymb}
\usepackage[colorlinks=true,linkcolor=blue,urlcolor=blue]{hyperref}
\setlength{\parindent}{0pt}
\setlength{\parskip}{0.7em}

\begin{document}

\title{Explainer: \emph{Inference on Welfare and Value Functionals under Optimal Treatment Assignment}}
\author{}
\date{}
\maketitle

This note gives an intuition-first explanation and then states the paper's formal distinction between welfare and value functionals.

\section{Feynman Approach for Intuition}

Suppose a planner treats a person if and only if that person's conditional average treatment effect, or CATE, is nonnegative. Then there are two natural summary questions.

The first is welfare: how much total gain do we get if we treat exactly those people with nonnegative CATE?

The second is value: what is the average of some other quantity over the people we would treat under that rule? For example, what share of the population gets treated, or what is the average of a covariate among treated people?

These sound similar, but the paper shows they are statistically very different. The reason is geometric. The treatment rule changes at the boundary where CATE equals zero. For welfare, the integrand itself is zero on that boundary, so boundary uncertainty drops out at first order. For a general value functional, the integrand usually does \emph{not} vanish there, so uncertainty about where the zero-CATE boundary lies matters directly.

\section{Formal Setup}

Let
\[
h_0(x)=\mathrm{CATE}(x)=\mathbb E[Y(1)-Y(0)\mid X=x].
\]
The optimal treatment rule is to treat type $x$ whenever $h_0(x)\ge 0$.

The welfare functional is
\[
W(h_0)=\int [h_0(x)]_+ f(x)\,dx,
\]
where $[t]_+=\max\{t,0\}$.

The more general value functional is
\[
V(h_0)=\int \mathbf 1\{h_0(x)\ge 0\}v_0(x)f(x)\,dx,
\]
where $v_0(x)$ is a user-chosen value object.

The level set
\[
\mathcal M_0=\{x:h_0(x)=0\}
\]
is the boundary between treated and untreated types under the optimal rule.

\section{The Key Analytical Difference}

The paper derives functional derivatives using a generalized Leibniz rule. This is where the whole result comes from.

For $W(h_0)$, perturbing $h_0$ changes both the integrand and the treatment region, but the boundary term vanishes because $h_0(x)=0$ on $\mathcal M_0$. So the first-order expansion is comparatively regular.

For $V(h_0)$, perturbing $h_0$ also moves the treatment region, and now the boundary term does \emph{not} vanish because the integrand is $v_0(x)f(x)$ rather than $h_0(x)f(x)$. Therefore the leading expansion contains a thin-set integral over $\mathcal M_0$.

That is the whole paper in one sentence: welfare is regular because the boundary contribution disappears, while value is typically irregular because the boundary contribution survives.

\section{Main Rigorous Results}

The paper proves four main inference statements.

\textbf{1. Welfare with known density.}
Theorem 1 establishes $\sqrt n$ asymptotic normality for plug-in estimators of the welfare functional under known target density.

\textbf{2. Welfare with unknown density.}
Theorem 2 extends the result when the density is not known and the integral is replaced by the empirical distribution of covariates. The same regularity logic survives.

\textbf{3. Value with known density.}
Theorem 3 shows that the value functional is generally not $\sqrt n$ estimable. Instead, its asymptotic distribution occurs at a one-dimensional nonparametric rate because the first-order term depends on the zero-CATE boundary.

\textbf{4. Value with unknown density.}
Theorem 4 shows the empirical-density version has the same asymptotic behavior as the known-density case.

\section{Why Sieve Inference Appears}

Once the leading term for the value functional is a thin-set integral, variance estimation is no longer a routine delta-method exercise. The paper therefore uses sieve methods and the sieve Riesz representer to estimate the asymptotic variance and build feasible confidence intervals.

This is not technical decoration. It is forced by the structure of the problem. The irregularity comes from the geometry of the unknown boundary, so the inferential machinery has to keep track of that geometry.

\section{How to Read the Paper}

The cleanest reading strategy is:
\begin{enumerate}
\item understand the treatment rule $h_0(x)\ge 0$;
\item focus on the boundary $\mathcal M_0=\{x:h_0(x)=0\}$;
\item compare what happens to the first-order derivative of $W$ and $V$ when that boundary moves;
\item interpret the different rates as a direct consequence of whether the boundary term vanishes.
\end{enumerate}

Everything else in the paper, including the Monte Carlo results and empirical application, is built on that derivative comparison.

\section{Bottom Line}

At the intuition level, welfare is easier because on the knife-edge types the gain from treatment is exactly zero.

At the rigorous level, that makes the welfare functional regular and root-$n$ estimable, while general value functionals inherit thin-set irregularity through a boundary integral over the zero-CATE set.

\end{document}
